\begin{helper}[Replacement-tree confinement]
Let $L$ be a \noderef{ScheideLastExitPath} certificate for $\alpha$, and
suppose its last inside vertex $v$ does not belong to
$X_0=\noderef{TreeSequenceVertexSet}(T)$.  Then there is a defective
color $\delta$ such that $v\in U_\delta=\{(R\delta).u1,(R\delta).u2\}$.
Write $r=D.\gammaIndex(\delta)$,
$\gamma=C.\gamma(r)$,
$Z=X_0\setminus\{C.\mathit{blockVertex}(r,0),
C.\mathit{blockVertex}(r,1)\}$,
$T_\delta=(R\delta).\mathit{replacementTree}$, and
\[
 \varphi_\delta=\noderef{TwoColorComponentSwap}
   (G,\varphi,\gamma,\delta,\operatorname{inl}D.v0(\delta)).
\]
The labeled $(D.\alpha_0,\alpha)$ incidence graph is the same under
$\varphi_\delta$ and $\varphi$.  Every color used on $T_\delta$ under
$\varphi_\delta$ is missing under $\varphi_\delta$ at a vertex of $Z$.
The vertex $v$ misses a member of
$\noderef{TreeFreeColors}(\varphi_\delta,T_\delta)$, and $\alpha$ is not
in \noderef{DefectiveColors}$(G,\varphi_\delta,V(T_\delta))$.
Every vertex-node of $L$'s path from $y_0$ through $v$ belongs to
$V(T_\delta)$.  Finally, every edge identity belonging both to
$\noderef{BoundaryEdges}(G,X)$ and to the full
$(D.\alpha_0,\alpha)$ component through $y_0$ is the crossing edge of
$L$.
\end{helper}
\begin{proof}
Since $v\in X\setminus X_0$, the indexed-union definition
\noderef{ScheideEnlargedVertexSet} supplies a defective index $\delta$
with $v\in U_\delta$.  The appropriate one of $(R\delta).f1$ and
$(R\delta).f2$ is a $\delta$-colored edge incident with $v$ and with its
other endpoint in $Z$, by \noderef{ScheideReplacementData}.  The second
conclusion of \noderef{ScheideSymmetricDefectiveInternal} supplies a
$\gamma$-colored edge internal to $U_\delta$, also incident with $v$.
Both are distinct from the crossing $g$ of $L$, because their endpoints
belong to $X$ whereas $g$ leaves $X$.  The extension field of
\noderef{ScheideReducedConfiguration} supplies the structural
configuration.  Properness, obtained by expanding
the \noderef{GlobalMaximumTriple} and \noderef{KTriple} inside
\noderef{ScheideStructuralConfiguration} to its
\noderef{IsProperEdgeColoring} field, and
\noderef{AdjacentEdges} therefore give
\[
 \alpha\ne\delta,\qquad \alpha\ne\gamma.                 \tag{1}
\]
The setup \noderef{ScheideDefectiveSetup} gives
$\alpha_0\ne\gamma$.  Defectiveness, expanded through
\noderef{DefectiveColors} and \noderef{DefectiveColor}, supplies a
$\delta$-colored boundary edge of $X_0$, while
\noderef{ScheideGlobalMaximumOrder} makes
$X_0$ closed and $y_0\in X_0$ misses $\alpha_0$.  Thus
\noderef{ColoredVertexSetClosed} implies
\[
 \alpha_0\ne\delta,\qquad \alpha_0\ne\gamma.             \tag{2}
\]

The formula in \noderef{TwoColorComponentSwap} changes only the values
$\gamma$ and $\delta$.  Equations (1)--(2), applied in both directions,
give for every edge identity $q$
\[
 \varphi_\delta(q)=\alpha_0\iff\varphi(q)=\alpha_0,
 \qquad
 \varphi_\delta(q)=\alpha\iff\varphi(q)=\alpha.           \tag{3}
\]
Expanding \noderef{TwoColorIncidenceGraph}, (3) proves equality of the
two labeled incidence graphs, and hence preserves the full component and
the path $P$ recorded by $L$.

The theorem \noderef{ScheideReplacementMaximum} says that
$(T_\delta,e_0,\varphi_\delta)$ is globally maximum.  Applying
\noderef{ScheideGlobalMaximumOrder} makes $V(T_\delta)$ elementary and
closed under $\varphi_\delta$.  The exact vertex-set field of
\noderef{ScheideReplacementData} is
\[
 V(T_\delta)=Z\cup U_\delta.                              \tag{4}
\]
The exact base and block positions in
\noderef{ScheideStructuralConfiguration}, together with vertex
injectivity in \noderef{TreeSequence}, place $y_0$ in $Z$.  The exit
witness fields of \noderef{ScheideDefectiveSetup}, unfolded through
\noderef{ScheideExitVertex}, place $D.v0(\delta)$ in
\noderef{ScheideExitVertices}; then
\noderef{ScheideStructuralExitsInBase} places it in the base and hence
in $Z$.  In particular $y_0,D.v0(\delta),v$ all lie in $V(T_\delta)$.

The missing-set formula in \noderef{ScheideDefectiveSwitchEffect}
shows that $\alpha_0$ remains missing at $y_0$.  Away from
$D.v0(\delta)$ the missing set is fixed.  At that vertex it becomes
$(\overline\varphi(D.v0(\delta))\setminus\{\gamma\})
\cup\{\delta\}$, which retains $\alpha_0$ by (2).

We prove the used-color assertion explicitly.  Let $c$ be in
\noderef{TreeUsedColors}$(\varphi_\delta,T_\delta)$.  The exact edge-set
and position fields of \noderef{ScheideReplacementData} say that its
witnessing edge is either a retained old tree edge or one of the appended
edges $(R\delta).f1,(R\delta).f2$.  The deleted block starts at position
$p-1\geq1$, so the old first edge remains the first replacement edge and
cannot witness membership in \noderef{TreeUsedColors}.  The base-position
and block-position fields of \noderef{ScheideStructuralConfiguration},
together with vertex injectivity in \noderef{TreeSequence}, also show
that every base vertex is in $Z$.

For a retained noninitial base edge, the
\noderef{TashkinovTree} field of the original globally maximum triple
supplies an earlier missing witness in the base.  Its color is not
$\gamma$ by the base-avoidance field of
\noderef{ScheideStructuralConfiguration}.  For a retained block edge,
the structural block-color and base-missing-witness fields give a base
witness, and injectivity of the block colors shows again that its color
is not the deleted color $\gamma$.  Such an edge is not $D.f0(\delta)$,
because every displayed old tree edge has both endpoints in $X_0$ by
\noderef{TreeSequence}, whereas $D.f0(\delta)$ has an endpoint outside
$X_0$.  The fixed-edge and missing-set clauses of
\noderef{ScheideDefectiveSwitchEffect} therefore preserve the edge color
and its witness; even if the witness equals $D.v0(\delta)$, the
erase-insert formula preserves every color other than $\gamma$.  Hence a
retained used color is missing under $\varphi_\delta$ at a vertex of $Z$.

For an appended edge, \noderef{ScheideReplacementData} gives old color
$\delta$, an endpoint in $Z$, and inequality from $D.f0(\delta)$.
The fixed-edge clause of \noderef{ScheideDefectiveSwitchEffect} leaves
its color equal to $\delta$, while that node's missing-set formula
inserts $\delta$ at $D.v0(\delta)\in Z$.  Thus in all cases
\[
 c\in\noderef{TreeUsedColors}(\varphi_\delta,T_\delta)
 \Longrightarrow
 \exists z\in Z,
 c\in\noderef{MissingColors}(G,\varphi_\delta,z).           \tag{5}
\]

Apply \noderef{ScheideGlobalMaximumMissingCounts} to the replacement
triple and $v\in V(T_\delta)$.  Its two alternatives both give a color
$\beta\in\noderef{MissingColors}(G,\varphi_\delta,v)$.  If $\beta$ were
used on $T_\delta$, (5) would give $z\in Z$ also missing $\beta$.
Because $z\in X_0$ and $v\notin X_0$, these are distinct vertices of
$V(T_\delta)$, contradicting its
\noderef{ColoredVertexSetElementary} property.  Thus $\beta$ is unused;
the definitions \noderef{TreeMissingColors} and
\noderef{TreeFreeColors} show that $\beta$ is a free color missing at
$v$.

If $\alpha$ were defective for $T_\delta$, the preserved full path $P$
from $y_0$, the missing color $\alpha_0$ there, and the fact that $v$ is
the last replacement-tree vertex of $P$ would satisfy
\noderef{ScheideExitVertex}.  Hence
$v\in\noderef{ScheideExitVertices}(T_\delta,\varphi_\delta)$.  But
\noderef{ScheideExitNoFree}, applied to the replacement globally maximum
triple, says that such a vertex misses no free color, contradicting the
chosen $\beta$.  Therefore $\alpha$ is not defective for
$T_\delta$.

Closedness of $V(T_\delta)$ and the missing color $\alpha_0$ at $y_0$
forbid every $\alpha_0$-colored boundary edge of that tree.  The edge
$g$ leaves $X$ from $v$, so by (4) it is an $\alpha$-colored boundary
edge of $V(T_\delta)$.  If $P$ left $V(T_\delta)$ before reaching $v$,
its first departure would also have color $\alpha$ and would be distinct
from $g$, contradicting the nondefectiveness just proved.  Hence every
vertex-node of $P$ through $v$ lies in $V(T_\delta)$.

Finally, the prefix just proved lies in $V(T_\delta)\subseteq X$, while
the last-index field of \noderef{ScheideLastExitPath} puts every later
vertex-node outside $X$.  Its support equality says that every edge-node
of the full component lies on $P$.  Therefore the only edge-node of that
component with one endpoint in $X$ and the other outside is the edge
immediately following $v$, namely $g$.  This proves the asserted final
uniqueness.
\end{proof}
